\ssszh 接下来把球与 Whittaker 两类范畴配对，得到
%
\[
   \Dmod(L^+G \backslash LG / LN, \psi) \simeq \QCoh( \pt / \cG).
\]
%
这同样是 Frenkel--Gaitsgory--Vilonen \cite{fgv} 的一条已知而且要重的定理。可以把它理解为 $p$-adic 表示论中
Casselman--Shalika 公式 \cite{cassshalika} 的几何精细化，该公式显式地
确定了一个非分歧表示的 Whittaker 余向量。

这里的参数配比类似。也就是说，$LG$ 被 double cosets
\[
   L^+G \cdot t^{\check{\mu}} \cdot LN, \quad \check{\mu} \in \check{\Lambda}
\]
分层，而一条轨道承载 Whittaker 层当且仅当其余权 $\check{\mu}$ 反支配。

\begin{re} 我们还想简短评注与双球情形（即导出 Satake）的一处基本差别。在那里，谱侧是
$\on{Rep}(\cG)$ 的一个导出增强。粗略地说，之所以如此，是因为
$L^+G \backslash LG / L^+G$ 上相交上同调层的茎沿着轨道闭包相当有趣，
对应着 $\cG$-表示的 $q$-权数重数 \cite{lusztigqchar}。

在当下情形，没有这样的导出增强。之所以如此，是因为
$L^+G \backslash LG / LN, \psi$ 上的相交上同调对象是 clean 的，
即在支撑中开轨道的边界处，其 $!$-茎与 $*$-茎都消没。

\end{re}

\ssszh \label{ss:sphericalkmmod}最后，把球向量与 Kac--Moody 表示配对，得到
\[
   \fgc\mod^{L^+G} \simeq \QCoh(\Op^{\on{unramified}}),
\]
其中后一个模空间按定义参数化平凡丛上的 oper 结构，即
\[
\Op^{\on{unramified}} \simeq \pt / \cG \underset \LS \times \Op.
\]

这同样是一条已知的定理，属于 Frenkel--Gaitsgory \cite{fg-sph}，它包含并
加强了先前关于临界水平 $L^+G$-等变 Kac--Moody 表示的许多结果，
这些结果属于 Feigin--Frenkel、Beilinson--Drinfeld 等人。

此时，Grothendieck 群不再是可数的。取而代之，单模构成由射有限
维仿射空间参数化的连续族。尽管如此，我们仍能给出如下粗略的
计数。

在自守侧，有一组紧生成元，即 Weyl 模；按定义，它们是有限维不可约
$G$-模的抛物诱导
\[
  \mathbb{V}_\lambda := \on{pind}_{\fg}^{\fgc}(L_\lambda), \quad \lambda \in \Lambda^+.
\]
在当下的李代数情形，抛物诱导由对应
\[
  \fg \leftarrow \fg[[t]] \rightarrow \fgc
\]
给出，也就是说，它是限制函子与诱导函子的复合
%
\[
\on{pind}_{\fg}^{\fgc}: \fg\mod \xrightarrow{\on{Res}} \fg[[t]]\mod \xrightarrow{\on{ind}} \fgc\mod.
\]
%

每个 $\mathbb{V}_\lambda$ 的自同态由中心 $Z(\fgc)$ 的一个商给出，并且它在
$\gkc\mod^{L^+G}_\lambda$ 中是投射的——这里 $\gkc\mod_\lambda$ 指 $Z(\fgc)$ 通过同一个商
作用的 $\gkc$-模范畴。这是一个非平凡的事实。

\begin{ex}\label{e:regop} 若 $\lambda = 0$，就得到临界水平处的真空
模\footnote{这里「真空」一词源于它在共形场论中的出处。也就是说，它是某个顶点算子代数的
底向量空间，而后者作为局部算子出现在具有 $G$-对称性的手征共形场论中。}
	\[
	  \mathbb{V}_0 \simeq \on{pind}_{\fg}^{\fgc}(k).
	\]
在 Feigin--Frenkel 同构之下，它的自同态典范地等同于非穿孔形式圆盘上
oper 代数上的函数代数
%
\[
    \Op^{\on{reg}} \simeq \{ d + f dt + \check{\fb}[[t]]dt \} / \cN[[t]].
\]
%
\end{ex}

在谱侧，随着 $\lambda$ 变化，这些自同态代数对应 $\Op$ 内部互不相交的闭
子概形。模空间 $\Op^{\on{unramified}}$ 是它们并集的形式加厚，而该定理进一步把这些
幂幺法方向与 Weyl 模的自扩张相匹配。

\ssszh\label{s:2realz} 可以进入图 \ref{f:tableaut} 与图 \ref{f:tablespec} 的下一行了，即
驯分歧与幂幺单值的情形。先遇到的是自同态，亦即一个幺半等价
%
\begin{equation} \label{e:2realz}
 \Dmod(I \backslash LG / I) \simeq \ICn( \widetilde{\sN} / \cG \underset {\cfg
 / \cG} \times \widetilde{\sN} / \cG).
\end{equation}
%
这一漂亮的陈述在多个意义上都值得注意。其一，有限 Hecke 范畴并没有这样的
凝聚实现，在这个意义上，该等价确实是仿射情形（具体地说，是其与算术的
联系）所特有的。\footnote{诚然，通过 Soergel 双模的构造可以把单值 Hecke
范畴写成某些凝聚层在
\[
\check{T} \underset {\check{T}/W_f}{\times} \check{T}
\]
上的同伦范畴，参见 \cite{bezrichetop}。不过这一等同在某种程度上更微妙，因为
它只在倾斜对象上、且仅在非导出意义下才全忠实；也就是说，相应的
Soergel 双模在双模范畴中有非平凡的 Ext。无论如何，\eqref{e:2realz} 还有另一个
令人愉快的特性：有限 Hecke 范畴的 Soergel 双模表述确实嵌入其中。}
如同我们迄今遇到的其它低分歧陈述，它是 Bezrukavnikov \cite{bez2realz} 的一条著名定理。

据我们所知，这一等价最初是作为猜想由 Ginzburg \cite{chrissginzburg} 提出的。一个主要
动机来自 $p$-adic 情形：在那里，$\on{Dist}(I \backslash G(\Fqt) / I)$（即仿射 Hecke 代数）
与 Steinberg 簇
\[
\widetilde{\sN} \underset \cfg \times \widetilde{\sN}
\]
的等变 $K$-理论之适当特化之间的类似同构，属于 Kazhdan--Lusztig \cite{kazhdanlusztig87}。他们以及
独立地 Ginzburg，用这一同构分类了仿射 Hecke 代数的不可约表示并确定了它们的特征标。等价地说，他们
确定了 $G(\Fqt)$ 的非分歧主系列的不可约子商的特征标。

来数参数。在自守侧，$LG$ 的相关分层其成分是
%
\[
    I\cdot w \cdot t^{\check{\mu}} \cdot I, \quad w \in W_f,\ \check{\mu} \in \check{\Lambda}.
\]
%
这是 $G$ 的 Bruhat 分解（双陪集 $B \cdot w \cdot B$，$w \in W_f$）的仿射类比，而这里的指标集
确实就是（扩张）仿射 Weyl 群 $W_f \ltimes \check{\Lambda}$。注意，如第 \ref{s:abg} 节那样，
我们再一次把 $L^+G$ 缩小到 $I$，因此分层多出了一个 $W_f$ 因子。

在谱侧，这些参数如下出现。首先，Steinberg 簇\footnote{虽然这是标准术语，我们要强调，
对于该幺半等价，我们实际用的是它的导出增强。}显式地参数化三元组
\[
(X, \cB', \cB''),
\]
其中 $X$ 是 $\cfg$ 的幂幺元，而 $\cB'$、$\cB''$ 是包含它的两个 Borel 子群。Steinberg 簇是
$\widetilde{\sN} \times \widetilde{\sN}$ 中 $W_f$ 个局部闭 Lagrange 子簇的并，其中每个
Lagrange 子簇由 $\cB_1$ 与 $\cB_2$ 的相对位置指定，即由 $\cB \backslash \cG / \cB$ 的一个点指定。

\begin{re} 回忆我们真正关心的是商
\[
\widetilde{\sN} / \cG \underset {\cfg/\cG} \times \widetilde{\sN} / \cG.
\]
把它认作 Hecke 叠
\begin{equation} \label{e:strestopesto} \cB \backslash \cG / \cB \simeq \Fl_{\cG} \times \Fl_{\cG} / \cG
\end{equation}
的余切（导出叠）会很有启发，这里的呈现方式是从 \eqref{e:strestopesto} 右端出发做 Hamilton 约化。
与此相关，所涉及的这些 $W_f$ 个成分不过是相应 Schubert 胞腔的余法丛。
\end{re}

如第 \ref{s:abg} 节那样，每个成分带有 $\check{\Lambda}$ 多个等变线丛，它们由沿任一投影
到 $\widetilde{\sN}/\cG$ 的拉回给出，\footnote{从 $\check{\Lambda}$ 到其 Picard 群的两个映射相差
$w$ 的相应元素的作用。}于是对成分求和，就得到所要的 $W_f \ltimes \check{\Lambda}$ 计数。

\ssszh 现在把非分歧主系列与 Whittaker 模型配对，得到等价
%
\begin{equation} \label{e:arkbez}
    \Dmod( I \backslash LG / LN, \psi) \simeq \QCoh( \widetilde{\sN} / \cG).
\end{equation}
%
这一断言同样是已知的，是 Arkhipov--Bezukravnikov \cite{ab} 的一条非常好的定理。他们的论证被广泛视为
本领域的一颗明珠。该结果在 Bezrukavnikov 对 \eqref{e:2realz} 的证明中起基本作用，
这与 $K$-理论中一条类似恒等式在 Kazhdan--Lusztig 把仿射 Hecke 代数等同于 Steinberg 簇
的等变 $K$-理论时所起的作用相平行。

这里的参数计数如下。在自守侧，过渡到半无限变体更方便，即
\begin{equation} \label{e:semiinfre}  \Dmod( LN \cdot L^+T \backslash LG / LN,
\psi).\end{equation}
在这一呈现下，按 $LG$ 的陪集
\[ (LN \cdot L^+T) \cdot w \cdot t^{\check{\mu}} \cdot LN,   \quad w
\in W_f,\ \check{\mu} \in \check{\Lambda} \]
分层，一条轨道承载 Whittaker 层当且仅当 $w$ 等于 $W_f$ 的最长元 $w_\circ$。特别地，相关轨道
的集合典范地由 $\check{\Lambda}$ 参数化。

在谱侧，有
%
\[
    \widetilde{\sN} / \cG = (\cG \overset \cB \times \check{\fn}) / \cG \simeq
    \check{\fn} / \cB,
\]
%
而 $\check{\Lambda}$ 再次对应于从 $\pt / \cB$ 拉回的等变线丛。

关于这一等价，有两点评注。

\begin{re} 在 $p$-adic 情形，相关的分布空间
\begin{equation} \label{e:antisph}\on{Dist}( I \backslash
G(\Fqt) / N(\Fqt), \psi)\end{equation}
是仿射 Hecke 代数的一个诱导表示。更准确地说，若以 $H$ 记（扩张）仿射 Hecke 代数，$H_f$ 记有限 Hecke
代数，$k_{\on{sgn}}$ 记其符号表示，则 \eqref{e:antisph} 就是所谓反球表示
\[ H \underset {H_f} \otimes k_{\on{sgn}}.\]
几何情形有类似的公式，即
\[ \Dmod(I \backslash LG / LN, \psi) \simeq \Dmod(I \backslash LG / I)
\underset {\Dmod(B\backslash G / B)} \otimes \Dmod(B \backslash G/ N, \psi),\]
参见 \cite{ahc}。
\end{re}

\begin{re} 我们在 \eqref{e:semiinfre} 中把 \eqref{e:arkbez} 的自守侧改写成了半无限形式。这是为了
稍稍简化相关轨道的组合学。不过，记住这一替代表述还有不那么琐碎的理由。也就是说，回忆范畴表示的
抛物诱导确实给出仿射旗簇的半无限版本，参见第 \ref{s:unramps} 节。与此相关，在
Arkhipov--Bezukravnikov 的证明中，起中心作用的正是半无限来源的对象，例如 Wakimoto 层。\end{re}

\ssszh \label{ss:iwahorikmmod}把非分歧主系列与 Kac--Moody 表示配对，得到等价
%
\[
 \fgc\mod^I \simeq \QCoh(\widetilde{\sN}/\cG \underset \LS \times \Op).
\]
%
这一断言是前面讨论的 Frenkel--Gaitsgory 球情形结果（见第 \ref{ss:sphericalkmmod} 节）的延伸。
他们先在 \cite{fg06loc} 中提出猜想，几年后又通过对仿射旗簇上临界水平局部化的工作得到它的一个版本，
其中还对中心特征作了进一步特化 \cite{fglocaff}。\footnote{严格地说，在 {\em loc. cit.} 中给出的是
Iwahori 单值变体。}
